\documentclass[11pt]{article}
\usepackage[margin=1in]{geometry}
\usepackage{amsmath,amssymb,amsthm}
\usepackage[colorlinks=true,urlcolor=blue,linkcolor=blue]{hyperref}
\newtheorem{theorem}{Theorem}
\newtheorem{lemma}[theorem]{Lemma}
\newtheorem{remark}[theorem]{Remark}
\newcommand{\one}{\boldsymbol 1}
\newcommand{\tr}{\operatorname{tr}}
\newcommand{\R}{\mathbb R}
\title{A graph-edge obstruction to budget-controlled sharp spectral projectors}
\author{Partitioning research session; frozen papers unchanged}
\date{November 15, 2026}
\begin{document}
\maketitle

\paragraph{Status.}
This note proves a negative theorem, not a new clustering algorithm.
The mechanism is the classical rank-one orthogonality catastrophe;
the proof below gives an elementary finite, unit-mass weighted-graph
realization. Historical priority for this particular formulation has
not been established. In particular, no novelty is claimed for
rank-one secular equations or Cauchy overlap matrices.

\paragraph{The precise big-if target.}
Let $L'=L+dbb^T$, $b=e_p-e_q$, $d>0$, on a connected graph with
ordinary unit vertex masses, and put $\tau=db^TL^\dagger b$.
For $0<\alpha<1$, the existing response result controls
\[
 C_\alpha=(L^\dagger)^{\alpha/2}(L'^\alpha-L^\alpha)
                   (L^\dagger)^{\alpha/2},\qquad
 \tr\log(I+C_\alpha)=\alpha\log(1+\tau).
\]
An attractive target would transfer this dimension-free budget to
the sharp low-energy projectors $P_t=\mathbf1_{(0,t)}(L)$,
$Q_t=\mathbf1_{(0,t)}(L')$, for example
\begin{equation}\label{eq:target}
 \int_0^\infty\|P_t-Q_t\|_F^2\,\frac{dt}{t}
       \le K\log(1+\tau),
\end{equation}
with a universal $K$, or any finite bound on the number of projector
singular values above a fixed tolerance depending only on $\tau$.
Both targets are false, even if the reference positive spectrum is
arbitrarily close to $1$ and the cutoff lies in a common gap with equal
projector ranks. Thus a small value of $\lambda_2(L)$ is not the issue.

\begin{theorem}\label{thm:nogo}
Fix $0<\gamma\le1/32$ and let $r\ge64$ be divisible by $16$, $n=r+1$,
and $\delta=1/(2n)$. There is a connected positively weighted complete
graph on $n$ unit-mass vertices and a single positive edge edit with
\[
 \frac{\gamma\delta}{1+\delta}\le\tau\le\gamma\delta.
\]
For every integer $3r/8\le k\le5r/8$ and every cutoff
\[
 1+\frac\delta r(k+2\gamma)<t<1+\frac\delta r(k+1),
\]
$P_t,Q_t$ both have rank $k$ and
\begin{equation}\label{eq:HS}
 \|P_t-Q_t\|_F^2\ge\frac{\gamma^2}{16}
                       \log(1+\lfloor r/16\rfloor).
\end{equation}
Moreover, put $Q=2^{16}$ and
\[
 s=\#\{\ell\in\mathbb Z_{\ge0}: 2Q^\ell\le r/16\}.
\]
At least $2s$ singular values of $P_t-Q_t$ are at least $\gamma/20$.
In particular this number is $\Omega(\log n)$ as $n\to\infty$ while
$\tau\to0$. Finally
\begin{equation}\label{eq:intlower}
 \int_0^\infty\|P_t-Q_t\|_F^2\,\frac{dt}{t}
 \ge\frac{\delta\gamma^2}{128}
                       \log(1+\lfloor r/16\rfloor).
\end{equation}
Consequently the ratio of this integral to $\log(1+\tau)$ tends to
infinity for fixed $\gamma$.
\end{theorem}

\begin{proof}
Choose an isometry $U:\R^r\to\one^\perp\subset\R^n$ with
$U^Tb=\sqrt{2/r}\,\one_r$. Such an isometry exists since
$\|b\|=\sqrt2$ and an orthogonal map can match any two vectors of
equal norm. Set
\[
 \Pi=I-\one\one^T/n,\quad
 H=U\operatorname{diag}(1/r,2/r,\ldots,1)U^T,\quad
 L=\Pi+\delta H,\quad d=\gamma\delta/2.
\]
Here $H\one=0$ and $\|H\|\le1$. For $i\ne j$,
$L_{ij}=-1/n+\delta H_{ij}\le-1/(2n)<0$.
Thus $L$ is exactly a positively weighted complete-graph Laplacian;
its diagonal is the required weighted degree because its row sums vanish.
Its positive eigenvalues are $\lambda_i=1+\delta i/r$.
Also
\[
 \tau=\frac{\gamma\delta}{r}\sum_{i=1}^r\lambda_i^{-1},
\]
which proves the leverage bounds. On the common positive support,
\[
 U^TL'U=I+\frac\delta r\big(D+\gamma\one_r\one_r^T\big),
 \qquad D=\operatorname{diag}(1,\ldots,r).
\]
Write the eigenvalues of the matrix in parentheses as
$\nu_1<\cdots<\nu_r$. Positive rank-one strict interlacing gives
$j<\nu_j<j+1$ for $j<r$ and $\nu_r>r$.

For $r/4\le j\le3r/4$, write $\nu_j=j+\theta_j$.
Its secular equation is
\begin{equation}\label{eq:secular}
 1+\gamma\left(-\frac1\theta+S_j(\theta)\right)=0,
 \qquad S_j(\theta)=\sum_{i\ne j}\frac1{i-j-\theta}.
\end{equation}
For $0<\theta\le1/4$, paired positive and negative terms contribute
at most
$\sum_{h\ge1}2\theta/(h^2-\theta^2)<1$ in absolute value.
The unpaired harmonic tail is bounded by $\log4+1$ on this range of
$j$ (enlarging to this bound also covers the end rounding).
In particular $|S_j(\theta)|\le8$.
The left side of~\eqref{eq:secular} is negative at $\theta=\gamma/2$
and positive at $\theta=2\gamma$.
Strict monotonicity between its poles therefore gives
\begin{equation}\label{eq:bulk}
 \gamma/2\le\theta_j\le2\gamma.
\end{equation}
Choose normalized new eigenvectors with coordinates in the old basis
\[
 V_{ij}=\frac{N_j}{i-j-\theta_j},\quad
 N_j=\left(\sum_{i=1}^r(i-j-\theta_j)^{-2}\right)^{-1/2}.
\]
Their signs are immaterial. For a bulk $j$, the off-pole square sum is
at most $8$, whereas $\theta_j^{-2}\ge256$.
Consequently
\begin{equation}\label{eq:norms}
 \frac\gamma{\sqrt8}\le N_j\le2\gamma,
 \qquad |V_{ij}|^2\ge\frac{\gamma^2}{8(j-i+1)^2}\quad(i<j).
\end{equation}

The cutoff intervals in the theorem lie above $\lambda_k$ and
$\mu_k=1+\delta\nu_k/r$ and below $\lambda_{k+1}$ and $\mu_{k+1}$;
the projectors therefore both have rank $k$. Put $q=\lfloor r/16\rfloor$.
All $j=k+1,\ldots,k+q$ are in the bulk. Thus
\begin{align*}
 \tr(P_t(I-Q_t))
 &=\sum_{i\le k,\,j>k}|V_{ij}|^2\\
 &\ge\frac{\gamma^2}{8}
       \sum_{a=0}^{q-1}\sum_{b=1}^q\frac1{(a+b+1)^2}
 \ge\frac{\gamma^2}{32}\sum_{h=1}^q\frac1h
 \ge\frac{\gamma^2}{32}\log(q+1).
\end{align*}
The penultimate inequality retains pairs $a+b=h$, $1\le h\le q$,
and uses $h/(h+1)^2\ge1/(4h)$.
Equal ranks give $\|P_t-Q_t\|_F^2=2\tr(P_t(I-Q_t))$,
proving~\eqref{eq:HS}.

For the singular-value assertion, use row indices $i=k-a$ and column
indices $j=k+b$, with $a,b\in[R_\ell,2R_\ell-1]$ and
$R_\ell=Q^\ell$, $0\le\ell<s$.
These are disjoint blocks of bulk indices. Change a common sign so
the resulting cross-overlap entries are $N_{k+b}/(a+b+\theta_{k+b})$.
Compress rows using the unit constant vector on each block, and columns
using the normalized vector proportional to $1/N_{k+b}$ on each block.
The resulting $s$ by $s$ matrix is $T=K\operatorname{diag}(c_\ell)$,
where~\eqref{eq:norms} gives
\[
 \gamma/\sqrt8\le c_\ell\le2\gamma,
 \quad K_{\ell h}=\frac1{\sqrt{R_\ell R_h}}
 \sum_{a=R_\ell}^{2R_\ell-1}\sum_{b=R_h}^{2R_h-1}
                \frac1{a+b+\theta_{k+b}}.
\]
The diagonal satisfies $K_{\ell\ell}\ge1/5$.
For $\ell\ne h$,
$|K_{\ell h}|\le\sqrt{\min(R_\ell,R_h)/\max(R_\ell,R_h)}
=Q^{-|\ell-h|/2}$.
Hence the off-diagonal matrix has operator norm at most
$2/(\sqrt Q-1)<1/30$, by its row and column sums.
It follows that $\sigma_{\min}(K)>1/6$ and
$\sigma_{\min}(T)>\gamma/(6\sqrt8)>\gamma/20$.
Singular-value min--max transfers these $s$ singular values from the
compression to the cross-overlap matrix.
For equal-rank projections these singular values are the sines of
principal angles, and each nonzero principal-angle sine occurs twice
among singular values of $P_t-Q_t$. This proves the claim.

There are at least $r/4$ disjoint cutoff intervals used above,
each of length $(1-2\gamma)\delta/r$ and lying below $1+\delta$.
Using~\eqref{eq:HS} and $\delta\le1/2$ gives~\eqref{eq:intlower}
with the stated conservative constant. Since
$\log(1+\tau)\le\tau\le\gamma\delta$, the ratio is at least
$\gamma\log(q+1)/128$ and diverges.
\end{proof}

\begin{theorem}[Unbounded projector rank at each fixed leverage]
\label{thm:fixedtau}
For every prescribed $\tau_*>0$ and every integer $h$, there is a
connected unit-mass weighted complete graph, one positive edge edit
of leverage \emph{exactly} $\tau_*$, and a cutoff in a common spectral
gap with equal projector ranks, for which at least $h$ singular values
of $P_t-Q_t$ are at least $1/600$.
The reference positive spectrum lies in $[1,1+1/(2n)]$.
Thus, even at a fixed arbitrarily small leverage, no finite
dimension-independent count bound exists at a tolerance below $1/600$.
\end{theorem}

\begin{proof}
Use the same graph $L$ as above, and set
\[
 a_r=\frac1r\sum_{i=1}^r(1+\delta i/r)^{-1},\qquad
 d=\frac{\tau_*}{2a_r},\qquad
 \gamma=\frac{2d}{\delta}=\frac{\tau_*}{\delta a_r}.
\]
Then $db^TL^\dagger b=\tau_*$ exactly, and $\gamma\ge1/32$ for
all sufficiently large $r$. Smallness of $\gamma$ was used only
to locate the interlacing roots and to bound their normalizations.
We now use estimates uniform for all $\gamma\ge1/32$.

Put $c=1/64$ and
$g_j(\theta)=\sum_{i=1}^r(i-j-\theta)^{-1}$.
For bulk $j$, the preceding harmonic-tail bound gives
$g_j(c)=-64+S_j(c)\le-56$, so
$1+\gamma g_j(c)<0$.
Changing indices $i\mapsto r+1-i$ gives the exact identity
\[
 g_j(1-c)=-g_{r-j}(c)\ge56,
\]
where $r-j$ is again a bulk index. Hence the root of the secular
equation satisfies $c\le\theta_j\le1-c$ for all such $\gamma$.
The two closest pole terms in its squared normalization sum are
at most $2c^{-2}=8192$; all remaining terms sum to less than $4$.
Therefore
\[
 1/100\le N_j\le1.
\]
Take $k=r/2$ and any
\[
 1+\frac\delta r(k+63/64)<t<1+\frac\delta r(k+1).
\]
These cutoffs again have equal ranks $k$. The identical dyadic
compression proof gives $T=K\operatorname{diag}(c_\ell)$, now
with $c_\ell\ge1/100$ and still $\sigma_{\min}(K)>1/6$.
It gives $2s$ singular values above $1/600$. Choosing $r$ sufficiently
large makes $2s\ge h$, completing the proof.
\end{proof}

\paragraph{Why the fractional response can still be tiny.}
The elementary inequality $dbb^T\preceq\tau L$ and operator
monotonicity of $x^\alpha$ imply
\[
 0\preceq C_\alpha\preceq((1+\tau)^\alpha-1)I.
\]
Thus the entire relative fractional response tends to zero in operator
norm for the constructed sequence, while the hard-projector differences
have logarithmically many nonnegligible singular directions. There is
no contradiction: the discontinuous spectral cutoff magnifies a very
dense band of eigenvalues. Also $L^\alpha$ has exactly the eigenvectors
and eigenvalue order of $L$; its hard projectors only reparameterize $t$.

\paragraph{Literature boundary.}
The rank-one orthogonality-catastrophe mechanism and logarithmic
Anderson integrals precede this note; see K\"uttler--Otte--Spitzer,
\emph{Anderson's Orthogonality Catastrophe for One-dimensional Systems},
\href{https://arxiv.org/abs/1301.4923}{arXiv:1301.4923}, and
Gebert, \emph{On an integral formula for Fredholm determinants related
to pairs of spectral projections},
\href{https://arxiv.org/abs/1705.02796}{arXiv:1705.02796}.
Gebert's exact determinant formula has common-resolvent boundary and
spectral-shift conditions; it is not itself a gap-free projector bound.
Pushnitski's scattering results give broader noncompactness phenomena:
\href{https://arxiv.org/abs/math/0702253}{arXiv:math/0702253}.
This note supplies its finite graph realization directly, without
invoking an infinite-volume transfer or claiming novelty for the
underlying phenomenon.

\end{document}
